\section{Related Work}

\subsection{Instruction Following of LLMs}

Instruction-following evaluation has moved from broad human preference scoring toward explicit requirement checking.
IFEval introduced a reproducible benchmark of verifiable constraints such as length, keyword, and format requirements \citep{zhou2023ifeval}.
FollowBench extends this direction with multi-level fine-grained constraints \citep{jiang2024followbench}, while InFoBench decomposes instructions into atomic requirements for more interpretable scoring \citep{qin2024infobench}.
LIFBench studies whether models can follow instructions reliably in long-context settings \citep{wu2025lifbench}.

Recent work further shows that realistic settings require domain or scenario-specific constraint design.
AgentIF evaluates long agentic instructions with tool, conditional, and example-implied constraints \citep{qi2025agentif}.
SciIF argues that scientific answer correctness and scientific constraint compliance are distinct capabilities, and uses evidence-based judging to audit whether required constraints are explicitly satisfied \citep{su2026sciif}.
The closest financial work is FIFE, which evaluates complex financial instruction following using 88 human-authored prompts and chainable verifiers \citep{matlin2025fife}.
FinIF follows this line but targets a different scale and use case: workflow-grounded financial deliverables with both a 300-item hard test set and a 1,064-example training pool.

\subsection{Financial Capability Evaluation}

Financial NLP benchmarks usually measure language understanding, knowledge, reasoning, question answering, or task quality rather than explicit instruction adherence.
FLUE provides a financial language understanding suite across five NLP tasks \citep{shah2022flue}.
FinanceBench evaluates open-book financial question answering over public company materials with evidence strings \citep{islam2023financebench}.
FinEval and CFinBench evaluate Chinese financial knowledge and practical financial ability at large scale \citep{guo2025fineval,nie2025cfinbench}.
FinBen aggregates financial tasks across information extraction, textual analysis, QA, generation, risk management, forecasting, and decision-making \citep{xie2024finben}.
BizFinBench emphasizes business-driven real-world financial applications and combines objective and subjective metrics \citep{lu2025bizfinbench}.

These benchmarks are valuable for measuring whether models know finance or can solve finance-oriented tasks.
However, an operational financial output must also satisfy delivery constraints: source placement, calculation traceability, threshold preservation, missing-evidence handling, and decision-boundary language.
FinIF therefore treats instruction adherence as a first-class capability rather than a side effect of answer quality.

\subsection{Automated and Auditable Evaluation}

FinIF uses LLM judges only where deterministic rules are insufficient.
This choice is motivated by prior work showing that LLM-as-a-judge can scale evaluation of open-ended outputs, while also requiring careful rubric design because judges can exhibit bias and inconsistent thresholds \citep{liu2023geval,zheng2023mtbench}.
Accordingly, FinIF separates deterministic checks from semantic judging and reports strict item-level pass rate in addition to aggregate constraint accuracy.
This design follows the auditability motivation in AgentIF and SciIF: a model should not receive credit for implicit compliance when the required evidence or boundary statement is absent from the output.
